% Build main.tex first. This letter imports its actual page and line labels.
\documentclass[11pt]{article}
\def\ReviewMode{0}
% EDIT THIS FILE FIRST. These values are shared by the manuscript and response.
% Replace every example name, affiliation, email and manuscript identifier.
% Use \& for a literal ampersand, \% for a percent sign and \_ for an underscore.
\newcommand{\ManuscriptTitle}{A clear question about ecological communities}
% ShortTitle appears in running headers; keep it brief.
\newcommand{\ShortTitle}{Ecological communities}
\newcommand{\ArticleType}{Research article}
% Journal and manuscript ID are used in the response letter, not a title-page box.
\newcommand{\TargetJournal}{Journal name}
\newcommand{\ManuscriptID}{Manuscript ID}
\newcommand{\RevisionDate}{4 October 2026}
% AuthorNames is plain text for PDF metadata and the supplement heading.
% Update AuthorBlock below as well; it is the formatted title-page author list.
\newcommand{\AuthorNames}{Jane Doe and Alex Smith}
\newcommand{\CorrespondingEmail}{jane.doe@example.org}
\newcommand{\Keywords}{community ecology; population dynamics; quantitative methods}

% Match each superscript to an affiliation. Add authors with the same pattern.
% The reading copy retains all identities; it does not anonymize a manuscript.
\newcommand{\AuthorBlock}{Jane Doe$^{1}$ and Alex Smith$^{2}$\\[.5em]
  \small $^{1}$Department of Ecology, Example University\\
  \small $^{2}$Department of Applied Mathematics, Example Institute\\
  \small Correspondence: \texttt{\CorrespondingEmail}}

\usepackage{ecology}
\addbibresource{references.bib}
% m- is a namespace: \cref{m-sec:methods} reads sec:methods from main.aux.
% Keep main.pdf beside response.pdf if you distribute clickable manuscript links.
\externaldocument[m-]{main}[main.pdf]
\usepackage[normalem]{ulem}
\fancyhead[R]{\small\sffamily Response to reviewers}
\setlength{\parindent}{0pt}
\setlength{\parskip}{.6em}
\newcounter{reviewcomment}
\newcommand{\ReviewerID}{E}
\renewcommand{\thereviewcomment}{\ReviewerID.\arabic{reviewcomment}}
\providecommand{\theHreviewcomment}{}
\renewcommand{\theHreviewcomment}{\ReviewerID.\arabic{reviewcomment}}
% Start each reviewer with a unique ID (E, 1, 2, ...); this resets their comments.
% Comment numbers are automatic; do not type them into the comment text.
\newcommand{\reviewer}[2]{%
  \renewcommand{\ReviewerID}{#1}\setcounter{reviewcomment}{0}\section*{#2}}
\newtcolorbox{reviewcommentbox}{enhanced,breakable,colback=white,colframe=EcoAccent,
  boxrule=0pt,borderline west={2pt}{0pt}{EcoAccent},arc=0mm,
  left=3mm,right=2mm,top=2mm,bottom=2mm}
\newenvironment{reviewcomment}{\refstepcounter{reviewcomment}\begin{reviewcommentbox}
  \textbf{Comment \thereviewcomment.}\par\smallskip}{\end{reviewcommentbox}}
\newenvironment{reply}{\par\noindent\textbf{Response.}\ }{\par}
\newtcolorbox{manuscriptexcerpt}{enhanced,breakable,colback=EcoPale,colframe=EcoPale,
  boxrule=0pt,arc=1mm,left=3mm,right=3mm,top=2mm,bottom=2mm}
\newcommand{\added}[1]{\textcolor{EcoChange}{#1}}
\newcommand{\deleted}[1]{\textcolor{gray}{\sout{#1}}}
% USAGE: \wherechanged{line:start}{line:end}; supply labels without the m- prefix.
% Place corresponding \linelabel commands around the changed passage in the
% manuscript. The helper handles both same-page and cross-page ranges.
% Rebuild main.tex after edits, then response.tex, and check against the final PDF.
\newcommand{\wherechanged}[2]{%
  \begingroup
  \edef\FirstChangePage{\getpagerefnumber{m-#1}}%
  \edef\LastChangePage{\getpagerefnumber{m-#2}}%
  \ifx\FirstChangePage\LastChangePage
    page~\pageref{m-#1}%
  \else
    pages~\pageref{m-#1}--\pageref{m-#2}%
  \fi, lines~\ref{m-#1}--\ref{m-#2}\endgroup}
\begin{document}
{\huge\bfseries Response to reviewers\par}\medskip
{\large\ManuscriptTitle\par}\medskip
{\small\ManuscriptID\quad\textbullet\quad\TargetJournal\quad\textbullet\quad\RevisionDate\par}
\bigskip
Thank the editor and reviewers briefly. State which files accompany the revision. Comments below are invented examples, not actual peer-review correspondence. Blue text marks additions and struck-through text marks deletions inside manuscript excerpts. Page and line references point to the review copy compiled from \texttt{main.tex}.

% REPLACE THE EXAMPLES BELOW with the actual comments and your point-by-point replies.
% Keep reviewcomment, reply and manuscriptexcerpt roles separate.
\reviewer{E}{Editor}
\begin{reviewcomment}\label{comment:editor}
Please clarify the biological question and the scope of inference.
\end{reviewcomment}
\begin{reply}
Describe the change directly and explain how it addresses the concern. The introductory example is located at \wherechanged{line:question-start}{line:question-end}.
\begin{manuscriptexcerpt}
The example question is how \deleted{populations change} \added{density-dependent growth approaches a carrying capacity}. The model is illustrative and is not evidence about a real population.
\end{manuscriptexcerpt}
\end{reply}

\reviewer{1}{Reviewer 1}
\begin{reviewcomment}\label{comment:assumptions}
The model's assumptions should be explicit.
\end{reviewcomment}
\begin{reply}
Point to the relevant section, explain the assumption and state what it limits. In this example, assumptions appear in \cref{m-sec:methods} on page~\pageref{m-sec:methods}. A mathematical result is demonstrated in \cref{m-prop:limit}. Cite additional references when they substantiate the response \citep{may1972}.
\end{reply}

\begin{reviewcomment}
Please supply the intermediate steps of the derivation.
\end{reviewcomment}
\begin{reply}
The supplementary derivation is in \cref{m-sec:supp-derivation}, including \cref{m-eq:supp-integral}. These references are generated from labels, rather than typed figure or equation numbers. A page/line range still needs checking against the exact PDF submitted to the journal.
\end{reply}

\reviewer{2}{Reviewer 2}
\begin{reviewcomment}
Please improve reproducibility reporting.
\end{reviewcomment}
\begin{reply}
Explain what was deposited, where it can be accessed and how it maps to the analyses. Do not promise public access when data restrictions apply. For simulation-based work, the model-description protocol of \citet{grimm2006} provides a useful example.
\end{reply}
\printbibliography[heading=bibintoc,title={References cited in this response}]
\end{document}
